% BOOK-FILE-ID: f-57fdd85a393e
\subsection{初項による違い}\label{節：初項による違い}

ここまでは式の形に注目してきたが,実は初項の選び方も解法を左右する.
次の漸化式で実験してみよう.
\[
  a_1=c,
  \quad
  a_{n+1}=\frac{a_n^2-1}{2}
\]
初項を記号 $c$ のままにすると,
\begin{align*}
  a_2&=\frac{c^2-1}{2},\\
  a_3&=\frac{c^4-2c^2-3}{8},\\
  a_4&=\frac{c^8-4c^6-2c^4+12c^2-55}{128}
\end{align*}
一般に$a_n$は$c$の$2^{n-1}$次の多項式となり,どうにも簡単な形にはなりそうにない.

特に $c=2$ と固定すると,
\[
  a_1=2,
  \quad
  a_2=\frac{3}{2},
  \quad
  a_3=\frac{5}{8},
  \quad
  a_4=-\frac{39}{128}
\]
となる.倍角の公式も対称式の恒等式も,そのままでは使えそうにない.
こういうときはグラフで動きを見てみるのがよい.
図では$f(x)=\frac{x^2-1}{2}$とし,\,$a_1=2$から曲線$y=f(x)$へ縦に進み,直線$y=x$へ横に進む.
縦の移動で次の値を計算し,横の移動でそれを次の入力にする.
初項を変えれば,同じ漸化式でも別の軌道をたどる.
\BookFigCobweb

しかし,\,$c=2$で簡単な形が見えなかったからといって,他の初項でも同じとは限らない.
例えば,
\[
  c=1+\sqrt{2},
  \quad
  c=1-\sqrt{2}
\]
なら定数\bookterm{sequence}{数列}$a_n=c$となる.
一般の$c$では閉じた\bookterm{general}{一般項}が見つからなくても,
\[
  a_n=f^{\,\circ(n-1)}(c),
  \quad
  f(x)=\frac{x^2-1}{2}
\]
と書くことはできる.\footnote{$f^{\,\circ m}$は$f$を$m$回合成することを表し,\,$f^{\,\circ0}(c)=c$とする.関数の値の$m$乗とは異なる.}
ただし,これは漸化式を書き換えただけで,一般項を求めたことにはならない.それでも,値の動きや\bookterm{limit}{極限},近似値を調べることはできる.
付録では\bookterm{matrix}{行列}や\bookterm{generating}{母関数}を使って考察を広げる.

基本対称式型も,初項次第ではすぐに一般項が分かる.
\[
  a_1=-1,\quad a_{n+1}=a_n^2-2
\]
なら$a_n=-1$となる.定数解は\bookterm{fixed}{不動点}の式
\[
  \alpha=\alpha^2-2
\]
の解$\alpha=-1,2$を初項に選べば得られる.

定数解$a_n=-1$も,\,\sectionref{節：対称式を使う}の$\alpha+\alpha^{-1}=a_1$による表示に含まれる.
その計算を複素数へ広げよう.\footnote{$\alpha^2-a_1\alpha+1=0$には複素数の範囲で必ず非零の解がある.\,$|a_1|<2$では$\alpha$が実数にならないため,実数の$\alpha$だけを使った先の例から一歩広げている.}
特に$-2\le a_1\le2$なら,\,$a_1=2\cos\theta$と書けるので,
$\alpha=e^{i\theta}$と選ぶと
\[
  a_n=\alpha^{2^{n-1}}+\alpha^{-2^{n-1}}=2\cos(2^{n-1}\theta)
\]
という\bookterm{trigonometric}{三角関数}の形になる.

$a_1=-1$では$\theta=\frac{2\pi}{3}$.倍角を$2\pi$ごとに読み替えると,\,$\frac{2\pi}{3}$と$\frac{4\pi}{3}$を行き来する.
どちらも$\cos$の値が$-\frac{1}{2}$なので$a_n=-1$となる.角度は動いているのに,横座標だけを見ると止まって見えるのだ.

\begin{bookpoint}
  \textbf{POINT}\par
  一般項が見つからないときは,不動点を初項に選んだ場合や,初項の範囲を分けた場合を調べる.
\end{bookpoint}
